\subsection{K(C) 上的乘法结构}

设 K 是一个交换环，A 是环 K 上的一个双代数（III, §11, No.~4, 第 585 页），其余积为 c，余单位为 \(\gamma\)。除非另有说明，张量积都是在 K 上取的。设 E 与 F 是（左）A-模。张量积 \(E \otimes F\) 赋有一个 \((A \otimes A)\)-模结构，它由公式

\[
(a \otimes b) (x \otimes y) = a x \otimes b y\tag{11}
\]

（对 \(a, b \in \mathrm{A}\)，\(x \in \mathrm{E}\)，\(y \in \mathrm{F}\)）刻画。利用同态 \(c: \mathrm{A} \to \mathrm{A} \otimes \mathrm{A}\)，我们从这个 \((\mathrm{A} \otimes \mathrm{A})\)-模导出一个 A-模 \(c_*(\mathrm{E} \otimes \mathrm{F})\)（II, §1, No.~13, 第 221 页）。更精确地说，设 \(a \in \mathrm{A}\)；若 \(c(a) = \sum_{i} a_{i} \otimes b_{i}\)，则我们有

\[
a (x \otimes y) = \sum_ {i} a _ {i} x \otimes b _ {i} y\tag{12}
\]

（对 \(x \in E\)，\(y \in F\)）成立。记号滥用起见，我们也把这样得到的 A-模记作 \(E \otimes F\)。由 c 的余结合性立即得出，K-模的典范同构

\[
\varphi : (\mathrm{E} \otimes \mathrm{F}) \otimes \mathrm{G} \longrightarrow \mathrm{E} \otimes (\mathrm{F} \otimes \mathrm{G})
\]

对一切 A-模 E、F、G 都是 A-线性的。同样，若双代数 A 是余交换的，则从 \(E \otimes F\) 到 \(F \otimes E\) 的典范同构是 A-线性的。最后，用 \(K_{\gamma}\) 表示以加群 K 为底群、外部合成律为 \((a, x) \mapsto \gamma(a)x\) 的 A-模。从 \(K \otimes E\)（相应地，\(E \otimes K\)）到 E 的典范同构是从 \(K_{\gamma} \otimes E\)（相应地，\(E \otimes K_{\gamma}\)）到 E 的 A-模同构。

命题 9. —— 设 K 是一个交换环，A 是 K 上余单位为 γ 的一个双代数，C 是 A-模的类的一个加性集，满足下列性质：

\begin{enumerate}
\def\labelenumi{(\roman{enumi})}
\item
  每个 \(\mathcal{C}\) 型的 A-模都是投射的（* 或更一般地，平坦的*）K-模。
\item
  若 A-模 E 与 F 都是 C 型的，则 A-模 \(E \otimes F\) 亦然。
\item
  上面定义的 A-模 \(\mathrm{K}_{\gamma}\) 是 \(\mathcal{C}\) 型的。
\end{enumerate}

则在加群 \(\mathrm{K}(\mathcal{C})\) 上存在唯一的环结构，其乘法满足

\[
[ \mathrm{E} ] _ {\mathcal {C}} [ \mathrm{F} ] _ {\mathcal {C}} = [ \mathrm{E} \otimes \mathrm{F} ] _ {\mathcal {C}}
\]

（对一切 C 型的 A-模 E 与 F）。\(\mathrm{K}(\mathcal{C})\) 的单位元是 \([\mathrm{K}_{\gamma}]_{\mathcal{C}}\)。若双代数 A 是余交换的，则环 \(\mathrm{K}(\mathcal{C})\) 是交换的。

赋以合成律 \((\mathrm{E},\mathrm{F})\mapsto\mathrm{cl}(\mathrm{E}\otimes\mathrm{F})\) 后，集 C 是一个幺半群，其单位元为 \(\mathrm{cl}(\mathrm{K}_{\gamma})\)。于是，\(Z^{(\mathcal{C})}\) 典范地具有一个环结构，其乘法由公式

\[
\operatorname{cl} (\mathrm{E}) \operatorname{cl} (\mathrm{F}) = \operatorname{cl} (\mathrm{E} \otimes \mathrm{F})\tag{13}
\]

及单位元 \(\operatorname{cl}(\mathrm{K}_{\gamma})\) 刻画（III, §2, No.~6, 第 446 页）。

给定一个 C 型的 A-模 F 和一个正合列

\[
0 \longrightarrow \mathrm{E} ^ {\prime} \longrightarrow \mathrm{E} \longrightarrow \mathrm{E} ^ {\prime \prime} \longrightarrow 0\tag{\( \mathcal{E} \)}
\]

（由 C 型的 A-模组成），序列

\[
0 \longrightarrow \mathrm{E} ^ {\prime} \otimes \mathrm{F} \longrightarrow \mathrm{E} \otimes \mathrm{F} \longrightarrow \mathrm{E} ^ {\prime \prime} \otimes \mathrm{F} \longrightarrow 0\tag{\(\left(\mathcal{E}\otimes \mathrm{F}\right)\)}
\]

（由 \((\mathcal{E})\) 导出）是正合的，因为 F 在 K 上是投射的（\(^{*}\) 或更一般地，平坦的\(_{*}\)）（II, §3, No.~6, 第 251 页，命题 5 及 II, §3, No.~7, 第 257 页，推论 6）。用 No.~2 的记号，我们则有

\[
r _ {\mathcal {E} \otimes \mathrm{F}} = r _ {\mathcal {E}} \operatorname{cl} (\mathrm{F}).\tag{14}
\]

由此得出，\(\mathbf{Z}^{(\mathcal{C})}\) 中由形如 \(r_{\mathcal{E}}\) 的元素生成的子群 R 是环 \(\mathbf{Z}^{(\mathcal{C})}\) 的一个右理想，且同样可证它是 \(\mathbf{Z}^{(\mathcal{C})}\) 的一个左理想。按定义，\(\mathrm{K}(\mathcal{C})\) 是商群 \(\mathbf{Z}^{(\mathcal{C})}/\mathrm{R}\)；因此在 \(\mathrm{K}(\mathcal{C})\) 上存在唯一的环结构，其乘法对一切 C 型的 A-模 E 与 F 满足 \([E]_{\mathcal{C}}[F]_{\mathcal{C}} = [E \otimes F]_{\mathcal{C}}\)。其单位元是 \([K_{\gamma}]\)。

当双代数 A 是余交换的时，幺半群 C 是交换的；由此得出环 \(\mathbf{Z}^{(\mathcal{C})}\) 与商环 \(\mathrm{K}(\mathcal{C})\) 都是交换的。

注. —— 仅在命题 9 的假设 (i) 与 (ii) 下，关于直和的格罗滕迪克群 \(\mathrm{K}^{\prime}(\mathcal{C})\)（VIII, 第 188 页）具有唯一的环结构，其乘法满足 \([E]_{\mathcal{C}}^{\prime}[F]_{\mathcal{C}}^{\prime} = [E \otimes F]_{\mathcal{C}}^{\prime}\)。其单位元是 \([K_{\gamma}]_{\mathcal{C}}^{\prime}\)。若双代数 A 是余交换的，则环 \(\mathrm{K}^{\prime}(\mathcal{C})\) 是交换的。证明与命题 9 类似，因为群 \(\mathrm{K}^{\prime}(\mathcal{C})\) 可与 \(\mathbf{Z}^{(\mathcal{C})}/\mathrm{R}^{\prime}\) 等同，其中 \(R^{\prime}\) 是 \(\mathbf{Z}^{(\mathcal{C})}\) 中由形如 \(\mathrm{cl}(\mathrm{E}^{\prime} \oplus \mathrm{E}^{\prime\prime}) - \mathrm{cl}(\mathrm{E}^{\prime}) - \mathrm{cl}(\mathrm{E}^{\prime\prime})\) 的元素生成的子群（同前所引）。

在命题 9 的假设 (i)、(ii) 与 (iii) 下，环 \(\mathrm{K}(\mathcal{C})\) 称为 C 的格罗滕迪克环。当 K 是一个体且 C 是在 K 上为有限维的 A-模的类所成之集时，这些条件特别是满足的。于是我们有下面的结果。

推论. —— 设 A 是交换域 K 上余单位为 \(\gamma\) 的一个双代数。则在加群 \(\mathrm{R}_{\mathrm{K}}(\mathrm{A})\) 上存在唯一的环结构，其乘法满足

\[
[ \mathrm{E} ] _ {\mathcal {C}} [ \mathrm{F} ] _ {\mathcal {C}} = [ \mathrm{E} \otimes_ {\mathrm{K}} \mathrm{F} ] _ {\mathcal {C}}
\]

（对一切在 K 上为有限维的 A-模 E 与 F）。\(\mathrm{R}_{\mathrm{K}}(\mathrm{A})\) 的单位元是 \([\mathrm{K}_{\gamma}]_{\mathcal{C}}\)。若双代数 A 是余交换的，则环 \(\mathrm{R}_{\mathrm{K}}(\mathrm{A})\) 是交换的。

例. —— 1) 设 K 是一个交换域。设 G 是一个群，K{[}G{]} 是群 G 的代数。我们把 G 与它在 K{[}G{]} 中的典范像等同（III, §2, No.~6, 第 446 页）。

我们赋予 K{[}G{]} 一个双代数结构，其余积 c 与余单位 \(\gamma\) 由

\[
c (g) = g \otimes g, \quad \gamma (g) = 1 \quad (g \in G).\tag{15}
\]

设 \(\mathrm{E}\) 与 \(\mathrm{F}\) 是 \(\mathrm{K}[\mathrm{G}]\)-模；由公式 (12)，\(\mathrm{K}[\mathrm{G}]\)-模在 \(\mathrm{E} \otimes \mathrm{F}\) 上的结构由

\[
g (x \otimes y) = g x \otimes g y \quad (g \in \mathrm{G}, x \in \mathrm{E}, y \in \mathrm{F}).\tag{16}
\]

K{[}G{]}-模 \(K_{\gamma}\) 是向量空间 K，赋以由 \(g\lambda = \lambda\)（对 \(g \in G\)，\(\lambda \in K\)）定义的 G 的作用。

环 \(\mathrm{R}_{\mathrm{K}}(\mathrm{K}[\mathrm{G}])\) 也记作 \(\mathrm{R}_{\mathrm{K}}(\mathrm{G})\)。它是交换的；其乘法由 {[}E{]} {[}F{]} = {[}E \(\otimes_{\mathrm{K}}\) F{]} 给出，而 \(\mathrm{R}_{\mathrm{K}}(\mathrm{G})\) 的单位元是 \([K_{\gamma}]\)。

*2) 设 \(\mathfrak{g}\) 是交换域 \(K\) 上的一个李代数，\(U(\mathfrak{g})\) 是它的包络代数；我们把 \(\mathfrak{g}\) 与它在 \(U(\mathfrak{g})\) 中的典范像等同（Lie, I, §2, No.~7, 第 39 页，推论 2）。我们赋予 \(U(\mathfrak{g})\) 一个双代数结构，其余积 \(c\) 与余单位 \(\gamma\) 由

\[
c (\xi) = \xi \otimes 1 + 1 \otimes \xi , \quad \gamma (\xi) = 0\tag{17}
\]

（对 \(\xi\in\mathfrak{g}\)）给出（Lie, II, §1, No.~4, 第 115 页）。

设 \(\mathrm{E}\) 与 \(\mathrm{F}\) 是 \(\mathrm{U}(\mathfrak{g})\)-模；由公式 (17)，\(\mathrm{U}(\mathfrak{g})\)-模在 \(\mathrm{E} \otimes \mathrm{F}\) 上的结构由

\[
\xi (x \otimes y) = \xi x \otimes y + x \otimes \xi y\tag{18}
\]

（对 \(\xi\in\mathfrak{g}, x\in\mathrm{E}\) 与 \(y\in\mathrm{F}\)）刻画。

格罗滕迪克环 \(\mathrm{R}_{\mathrm{K}}(\mathrm{U}(\mathfrak{g}))\) 在 Lie, VIII, §7, No.~6, 第 139 页中也记作 \(\mathcal{R}(\mathfrak{g})\)。*

设 A 是一个交换环。我们可以把 A 视为其自身上的一个余交换双代数，其余积为从 A 到 \(A \otimes_{A} A\) 的自然同构，余单位为 \(Id_{A}\)（III, §11, No.~4, 第 585 页）。由命题 9，我们得到下面的结果。

命题 10. —— 设 A 是一个交换环，C 是 A-模的类的一个加性集，满足下列三个条件：

\begin{enumerate}
\def\labelenumi{(\roman{enumi})}
\item
  每个 \(\mathcal{C}\) 型的 A-模都是投射的（* 或更一般地，平坦的*）
\item
  若 \(\mathrm{E}\) 与 \(\mathrm{F}\) 是 \(\mathcal{C}\) 型的 A-模，则 A-模 \(\mathrm{E} \otimes_{\mathrm{A}} \mathrm{F}\) 也是 \(\mathcal{C}\) 型的。
\item
  A-模 A 是 \(\mathcal{C}\) 型的。
\end{enumerate}

则在加群 \(\mathrm{K}(\mathcal{C})\) 上存在唯一的环结构，对每一对 C 型的 A-模 E, F 满足 \([E]_{\mathcal{C}}[F]_{\mathcal{C}} = [E \otimes_{A} F]_{\mathcal{C}}\)。\(\mathrm{K}(\mathcal{C})\) 的单位元是 \([A]_{\mathcal{C}}\)。

